\begin{itemize}
\item \textbf{Elektroizolační kapaliny:} proud přenášejí ionty. \emph{Vlastní vodivost} vzniká disociací molekul samotné kapaliny, \emph{nevlastní} disociací nečistot a příměsí, např. vody. Pohyb nabitých koloidních částic způsobuje elektroforetickou vodivost.
\item \textbf{Iontové krystaly:} vodivost umožňují poruchy mřížky a přeskoky iontů mezi polohami. Schottkyho porucha vytváří vakance, Frenkelova vakanci a iont v mezimřížkové poloze.
\item \textbf{Skelné amorfní látky:} vodivost je iontová; v nepravidelné síti se pohybují zejména slabě vázané alkalické ionty Na$^+$ a K$^+$. Konduktivita závisí na jejich obsahu a pohyblivosti.
\end{itemize}
